\def\ChapterCode{DOK}
\chapter
[\CO{[DOK]} Dokumentation]
{Dokumentation}
\label{chp:DOK}

\ChapterIntro
{%
Die medizinische Dokumentation
muss eine Geschichte erzählen. Eine Geschichte die ein
Dritter, unbeteiligter, verstehen \E{muss}. Dieser
unbeteiligte Dritte soll allein anhand der
Schilderung in Ihrer Dokumentation fähig sein eine
Verdachtsdiagnose zu stellen. Weiters sollte er
nachvollziehen können, was mit dem Patient passiert ist.
Auch die Ereignisse davor oder die Einsatzsituation soll
für denjenigen verständlich sein (Unfallhergang,
Wohnsituation bei psychiatrischen Patienten, \ldots)
}
{%
\begin{description}
  \item[Maintainer:] --
  \item[Autoren:] Diverse
  \item[Reviewer:] Standard-Reviewprozess
  \item[Version:] {\VarDocVersionTypeName} ({\VarDocVersionTypeDescription})
  \item[SHA1:] {\DocGitCommit}
\end{description}

{\TxTeilkapitelAASS}
}








% 
% \dictum
% []
% {}
% 
% \dictum
% []
% {}
% 
% \dictum
% []
% {}
% 
% \dictum
% []
% {}
% 
% \dictum
% []
% {}
% 
% \dictum
% []
% {}

\clearpage % TEMP temp clearpage

\section{\HeadingKeyword{Allgemeines} zur Dokumentation}

\index{Dokumentation!medizinische Aspekte}
\label{topic:dokumentation-medizinisch}

\TextSlide{\TxQuerverweise}
{%
\begin{itemize}
  \item Rechtliche Aspekte:
  \FullPrettyRef{topic:dokumentation-rechtlich},
  \FullPrettyRef{topic:dokumentationspflicht-fallbeispiel}
  \item Dokumentation anderer Gesundheitsberufe (Anamnese):
  \FullPrettyRef{topic:dokumentation-andere-gesundheitsberufe}
\end{itemize}
}{%
\begin{itemize}
  \item Rechtlich:
  \prettyref{topic:dokumentation-rechtlich},
  \prettyref{topic:dokumentationspflicht-fallbeispiel}
  \item Andere Gesundheitsberufe:
  \prettyref{topic:dokumentation-andere-gesundheitsberufe}
\end{itemize}
}{}{}{}

\Text{Dokumentation -- Oder: Damit sich auch ein
Anderer auskennt \ldots} {Die medizinische Dokumentation
muss eine Geschichte erzählen. Eine Geschichte die ein
Dritter, unbeteiligter, verstehen \E{muss}. Dieser
unbeteiligte Dritte soll allein anhand der
Schilderung in der Dokumentation fähig sein eine
Verdachtsdiagnose zu stellen. Weiters sollte er
nachvollziehen können, was mit dem Patient passiert ist.
Auch die Ereignisse davor oder die Einsatzsituation soll
für denjenigen verständlich sein (Unfallhergang,
Wohnsituation bei psychiatrischen Patienten, \ldots).
}

\begin{Synopsis}
  \SynopsisItem{Die Dokumentation dient sowohl der rechtlichen
  Absicherung, als auch der Beschreibung der Situation
  gegenüber Dritter, welche am Einsatz nicht
  beteiligt waren.}
\end{Synopsis}


\Text
{Dokumentationspflicht}
{%
Schon nach dem Sanitätergesetz besteht für jeden Sanitäter
die Pflicht, die von ihm gesetzten
sanitätsdienstlichen Maßnahmen zu
dokumentieren.\ZFootnote{{\S}\,5 SanG} Darüber hinaus müssen
analog zu anderen Gesundheitsberufen auch alle anderen
medizinisch relevanten Sachverhalte schriftlich und
nachvollziehbar fixiert werden.

Hintergrund dieser Dokumentationspflicht ist der
\EE{Heilbehandlungsvertrag}.
Im österreichischen Recht hat grundsätzlich der Kläger den
Schaden, die Rechtswidrigkeit und das Verschulden des
Beklagten zu beweisen. Besteht aber zwischen Kläger und
Beklagten ein Vertrag, so wird gem. \Z{{\S}\,1298} ABGB das
Verschulden des Beklagten angenommen bzw. der Beklagte muss
beweisen, dass ihn kein Verschulden trifft
(\I{Beweislastumkehr}).
Um nun für den Beklagten die Möglichkeit zu schaffen, sein
Unverschulden zu beweisen, wurde die
\I{Dokumentationspflicht} eingeführt, mit der bei
ausreichender Dokumentation wieder der Kläger das
Verschulden des Beklagten zu beweisen hat.


Inhalt der Dokumentation sollte nicht nur die 
gesetzten Maßnahmen sein, sondern auch 
\begin{itemize}
  \item situationsrelevante
  Umstände (beim Patient; am BO; während des Transportes;
  {\dots}),
  \item  u.\,U. relevante Inhalte einer Aufklärung,
  \item Einwilligung bzw. Ablehnung durch den Patienten
  (Revers) sowie
  \item sonstige für den Sanitäter entscheidend wirkende
  Umstände.
\end{itemize}



Es muss \E{gründlich}, \E{ausführlich} und
\E{nachvollziehbar} dokumentiert werden.
Die Dokumentation soll vorrangig am Einsatzprotokoll selbst
oder, wenn nötig, auf einem beigelegten
Blatt Papier erfolgen.
}
{}
{}
{}
{}


\Text{NACA-Score}
{%
\label{topic:naca} %
\Z{Der \I{NACA-Score}\ZFootnote{NACA \Lang{Abkz.}: National
Advisory Committee for Aeronautics} diente ursprünglich der
Bewertung von Erkrankungen / Verletzungen als
Entscheidungshilfe welche Soldaten zuerst auszufliegen
wären.} Heutzutage dient der NACA-Score der groben
statistischen Erfassung.
Man unterscheidet die Stufen 1 bis 7:


\begin{table}[H]
\TableBaseModification
\Captionof{table}{NACA-Score mit beispielhafter Verteilung in einem Rettungsdienstbereich \Z{(bezogen auf Wien im Jahr
2009, nach  \citep{KontrollamtWienKA-K-13-09})}}

\begin{tabular}{clrr}\toprule\addlinespace
\noindent\TabHeading{NACA}
&
\noindent\TabHeading{Bedeutung}
&
\noindent\TabHeading{Anteil Pat.}
&
\noindent\TabHeading{Anteil \PC}
\\\addlinespace\midrule\addlinespace
 I & Keine ärztliche Therapie erforderlich & 10\,217 & 8,7
\\\addlinespace
 II & Ambulante Abklärung & 62\,207 & 52,8
\\\addlinespace
 III & Stationäre Abklärung & 38\,307 & 32,5
\\\addlinespace
 IV & Akute Lebensgefahr nicht ausgeschlossen & 2\,830 & 2,4
\\\addlinespace
 V & Akute Lebensgefahr & 1\,022 & 0,9
\\\addlinespace
 VI & Reanimation & 34 & 0,03	
\\\addlinespace
 VII & Tod & 312 & 0,3
 \\\addlinespace\midrule\addlinespace
 &&117\,849 & 100
\\\addlinespace\bottomrule
\end{tabular}
\end{table}
}{}{}{}{}


% \clearpage % TEMP temp clearpage
\section{\HeadingKeyword{Beispiele}}

Die hier angeführten Beispiele dienen der Demonstration, wie
ausführlich man dokumentieren könnte. Sie können vom sonst
regional üblichen und durch Formulare oder
Computerprogramme vorgegebenen Form abweichen.


\Layout{57}
{}
{}
{}
{%
\renewcommand{\FontTypewriter}{\ttfamily}
\begin{tabularx}{\linewidth}{XXXXXX}
\toprule\addlinespace
\multicolumn{6}{c}{\TabHeading{Einsatzprotokoll}}
\\\addlinespace\midrule\addlinespace
Berufung: 	& {\FontTypewriter 12-B-01}
&
Text:		& \multicolumn{3}{p{.45\linewidth}}{{\FontTypewriter Sturz, möglicherweise gefährliche Körperregion; Radfahrer; Zufahrt über Steinspornbrücke}}
\\\midrule
Art:		& {\FontTypewriter RTW}		& QU 	& {\FontTypewriter 01.04.2007 12:00}	& BO 	& {\FontTypewriter 12:05}
\\
ZA			& {\FontTypewriter 12:31}	& AO	& {\FontTypewriter 12:55}			& EB 	& {\FontTypewriter 13:02}
\\\midrule
Name: 		& {\FontTypewriter MUSTERMANN, Hannes}
&
Alter:		& {\FontTypewriter 36}
&
Geschlecht:	& {\FontTypewriter m}
\\\addlinespace\midrule
\multicolumn{6}{c}{\TabSub{Diagnosen}}
\\
\multicolumn{6}{p{.98\linewidth}}{{\FontTypewriter RQW Stirn, Zustand nach Sturz mit Fahrrad}}
\\\midrule
\multicolumn{6}{c}{\TabSub{Befunde}}
\\
Vitalwerte			& {\FontTypewriter 12:10} & RR - HF	& {\FontTypewriter 130/70}\MmHg* - {\FontTypewriter 90}\PerMin*		& Sp\ce{O2}	& {\FontTypewriter 99}\PC*
\\
Vitalwerte			& {\FontTypewriter 12:43} & RR - HF	& {\FontTypewriter 130/80}\MmHg* - {\FontTypewriter 81}\PerMin*	& Sp\ce{O2}	& {\FontTypewriter 99}\PC*
\\
BZ					& {\FontTypewriter xxx}\MgDl*& Temp.	& {\FontTypewriter xx}\,°C									& Schmerzen: 	&
\\
Weitere Befunde:	& \multicolumn{5}{p{.8\linewidth}}{{\FontTypewriter --- }}
\\
NACA:				& 2
&
GCS:				& 15
\\\addlinespace\midrule
\multicolumn{6}{c}{\TabSub{Anamnese, Bemerkungen}}
\\
\multicolumn{6}{p{.98\linewidth}}{{\FontTypewriter ---}}
\\\midrule
\multicolumn{6}{c}{\TabSub{Maßnahmen}}
\\
\multicolumn{6}{p{.98\linewidth}}{{\FontTypewriter ad KH}}
\\\midrule
\multicolumn{6}{c}{\TabSub{Ziel}}
\\
Zielabt.:		& {\FontTypewriter Unfallchirurgie}
&
Trpt. m. Arzt: 	& {\FontTypewriter nein}
&
Zustand:		& {\FontTypewriter gleich}
\\\addlinespace\bottomrule\addlinespace
\end{tabularx}
}
{Ein Einsatzprotokoll: Schlechte Dokumentation. Dieses Beispiel zeigt eine mangelhafte Dokumentation.
Anhand der Niederschrift ist nicht nachvollziehbar, was
vorgefallen ist. }
{}








\Layout{57}
{}
{}
{}
{%
\renewcommand{\FontTypewriter}{\ttfamily}
\begin{tabularx}{\linewidth}{XXXXXX}
\toprule\addlinespace
\multicolumn{6}{c}{\TabHeading{Einsatzprotokoll}}
\\\addlinespace\midrule\addlinespace
Berufung: 	& {\FontTypewriter 12-B-01}
&
Text:		& \multicolumn{3}{p{.45\linewidth}}{{\FontTypewriter Sturz, möglicherweise gefährliche Körperregion; Radfahrer; Zufahrt über Steinspornbrücke}}
\\\midrule
Art:		& {\FontTypewriter RTW}		& QU 	& {\FontTypewriter 01.04.2007 12:00}	& BO 	& {\FontTypewriter 12:05}
\\
ZA			& {\FontTypewriter 12:31}	& AO	& {\FontTypewriter 12:55}			& EB 	& {\FontTypewriter 13:02}
\\\midrule
Name: 		& {\FontTypewriter MUSTERMANN, Hannes}
&
Alter:		& {\FontTypewriter 36}
&
Geschlecht:	& {\FontTypewriter m}
\\\addlinespace\midrule
\multicolumn{6}{c}{\TabSub{Diagnosen}}
\\
\multicolumn{6}{p{.98\linewidth}}{{\FontTypewriter RQW Stirn, Zustand nach Sturz mit Fahrrad}}
\\\midrule
\multicolumn{6}{c}{\TabSub{Befunde}}
\\
Vitalwerte			& {\FontTypewriter 12:10} & RR - HF	& {\FontTypewriter 130/70}\MmHg* - {\FontTypewriter 90}\PerMin*		& Sp\ce{O2}	& {\FontTypewriter 99}\PC*
\\
Vitalwerte			& {\FontTypewriter 12:43} & RR - HF	& {\FontTypewriter 130/80}\MmHg* - {\FontTypewriter 81}\PerMin*	& Sp\ce{O2}	& {\FontTypewriter 99}\PC*
\\
BZ					& {\FontTypewriter xxx}\MgDl*& Temp.	& {\FontTypewriter xx}\,°C									& Schmerzen: 	&
\\
Weitere Befunde:	& \multicolumn{5}{p{.8\linewidth}}{{\FontTypewriter 
ABC unauff. 
\underline{Traumacheck:} RQW Stirn,
ca 50\,ml Blutverlust sonst unauffällig. \underline{Neuro:} grob
neurologisch unauffällig, Bew. klar, Pat zeitl, örtl, zur
Situation \& Person orientiert, Pupillen mittelweit, Reakt.
prompt und seitengleich, Kraft\&Gefühl seitengl. in OE und
UE. 
}}
\\
NACA:				& 2
&
GCS:				& 15
\\\addlinespace\midrule
\multicolumn{6}{c}{\TabSub{Anamnese, Bemerkungen}}
\\
\multicolumn{6}{p{.98\linewidth}}{{\FontTypewriter 
\underline{Anamn.:} \underline{Sympt:} Kopfschmerz seit
Aufprall. \underline{Krankh.:} keine bek. \underline{All.}
keine bek. \underline{Med.:} Aspirin heute früh. letzte
Tetanus-Impf. nicht bek. \underline{Ereignisse} Pat.
fuhr mit Fahrrad bei mittlerer Geschwindigkeit und prallte mit Stirn gegen Ast, kein
Sturz. Lt. Begleiter keine B.losigkeit, Unfallhergang
erinnerlich. 
\underline{Letzte} Mahlzeit: heute Früh 2 Semmeln

}}
\\\midrule
\multicolumn{6}{c}{\TabSub{Maßnahmen}}
\\
\multicolumn{6}{p{.98\linewidth}}{{\FontTypewriter 
Steriler Verband, nüchtern belassen,
Transport liegend mit erh. OK.
}}
\\\midrule
\multicolumn{6}{c}{\TabSub{Ziel}}
\\
Zielabt.:		& {\FontTypewriter Unfallchirurgie}
&
Trpt. m. Arzt: 	& {\FontTypewriter nein}
&
Zustand:		& {\FontTypewriter gleich}
\\\addlinespace\bottomrule\addlinespace
\end{tabularx}
}
{Ein Einsatzprotokoll: Bessere Dokumentation}
{}










\Layout{57}
{}
{}
{}
{%
\renewcommand{\FontTypewriter}{\ttfamily}
\begin{tabularx}{\linewidth}{XXXXXX}
\toprule\addlinespace
\multicolumn{6}{c}{\TabHeading{Einsatzprotokoll}}
\\\addlinespace\midrule\addlinespace
Berufung: 	& {\FontTypewriter 25-B-04}
&
Text:		& \multicolumn{3}{p{.45\linewidth}}{{\FontTypewriter Psychiatr. Erkr., unbekannter Zustand, Anrufer 3. Hand; Polizei vor Ort}}
\\\midrule
Art:		& {\FontTypewriter RTW}		& QU 	& {\FontTypewriter 01.05.2008 09:10}	& BO 	& {\FontTypewriter 09:20} 
\\
ZA			& {\FontTypewriter 10:00}	& AO	& {\FontTypewriter 10:23}				& EB 	& {\FontTypewriter 10:40}
\\\midrule
Name: 		& {\FontTypewriter MUSTERMANN, Max}
&
Alter:		& {\FontTypewriter 32}
&
Geschlecht:	& {\FontTypewriter m}
\\\addlinespace\midrule
\multicolumn{6}{c}{\TabSub{Diagnosen}}
\\
\multicolumn{6}{p{.98\linewidth}}{{\FontTypewriter Psychiatrische Erkrankung, V.a. Wahnvorstellungen, Gefährdung von Dritten, Unterbringung nach UbG}}
\\\midrule
\multicolumn{6}{c}{\TabSub{Befunde}}
\\
Vitalwerte			& {\FontTypewriter 09:30} 	& RR - HF	& {\FontTypewriter 130/70}\MmHg* - {\FontTypewriter 100}\PerMin*	& Sp\ce{O2}		& {\FontTypewriter 98}\PC* 
% \\
% Vitalwerte			& {\FontTypewriter 09:50} & RR - HF	& {\FontTypewriter 130/80}\MmHg* - {\FontTypewriter 100}\PerMin*	& Sp\ce{O2}		& {\FontTypewriter 99}\PC* 
\\
BZ					& {\FontTypewriter xxx}\MgDl*& Temp.	& {\FontTypewriter xx}\,°C 									& Schmerzen: 	&
\\
Weitere Befunde:	& \multicolumn{5}{p{.8\linewidth}}{{\FontTypewriter 
Pat. wach, agitiert, ABC soweit erhebbar unauff., \underline{Traumacheck:} keine äußeren Anz f Verletzungen,
sonst nicht erhebbar (unkoop.). \underline{Neuro:} Pat.
ist wach, sonst nicht erhebbar 

}}
\\
NACA:				& 3	
& 
GCS:				& 15
\\\addlinespace\midrule
\multicolumn{6}{c}{\TabSub{Anamnese, Bemerkungen}}
\\
\multicolumn{6}{p{.98\linewidth}}{{\FontTypewriter 
Pat.
sitzend in Wohnung angetr., von Pol. mit Handschellen gefesselt. Pat. warf lt. Polizei Gegstände aus dem Fenster
auf Straße, verletzte Passanten (ad RTW FLO-1). Pat.
gibt an, sich geg. Geheimdienst gewehrt zu haben,
jetzt agitiert und unkooperativ. 
\underline{Sympt:} keine Antwort
\underline{Krankh.:} lt. Spitalbrief Otto-Wagner-Sp. v.
1.4.08 paranoide Schizophrenie 
\underline{All.} k.A.
\underline{Med.:} nicht eruierbar, auf Tisch mehrere
angebr. Packungen Rohypnol 
\underline{Ereignisse} s.o.
\underline{Letzter Spitalsaufenthalt} k.A.
\underline{Wohnsituation:} heruntergekommen, vermüllt, Pat.
wohnt mit großem Hund, ganze Wohnung verkotet, bestialisch
stinkend.
}}
\\\midrule
\multicolumn{6}{c}{\TabSub{Maßnahmen}}
\\
\multicolumn{6}{p{.98\linewidth}}{{\FontTypewriter 
\underline{Maßnahmen:} Einweisung nach UbG, Transport in
Handschellen in Pol-begl. (Viktor 6, DNr. 47012). Protokoll
Amtsarzt ad. Spital. Wohnungsschlüssel ad Polizei (Viktor 2, DNr. 47011)}}
\\\midrule
\multicolumn{6}{c}{\TabSub{Ziel}}
\\
Zielabt.:		& {\FontTypewriter Psychiatrie}
&
Trpt. m. Arzt: 	& {\FontTypewriter n}
&
Zustand:		& {\FontTypewriter gleich}
\\\addlinespace\bottomrule\addlinespace
\end{tabularx}
}
{Einsatzprotokoll: Psychiatrischer und psychosozialer Patient. Hier ist nicht nur nachvollziehbar was passiert ist, auch
ist erkennbar dass versucht wurde, mögliche gefährliche
Diagnosen auszuschliessen (Schädel-Hirn-Trauma,
Hirnblutungen, \ldots). Bei diesem Patienten ist die Schilderung der Lebensumstände
wichtig, sie hilft der weiterbehandelnden Stelle beim
Entlassungmenagement -- der Patient kann ja schließlich
nicht in eine bestialisch stinkende, verkotete Wohnung
zurückgeschickt werden.}
{}

% \Layout{57}
% {}
% {}
% {}
% {%
% \renewcommand{\FontTypewriter}{\ttfamily}
% \begin{tabularx}{\linewidth}{XXXXXX}
% \toprule\addlinespace
% \multicolumn{6}{c}{\TabHeading{Einsatzprotokoll}}
% \\\addlinespace\midrule\addlinespace
% Berufung: 	& {\FontTypewriter MM-X-mm}
% &
% Text:		& {\FontTypewriter Berufungsdiagnose im Klartext, 3. Hand}
% &
% Zusatz:		& {\FontTypewriter }
% \\\midrule
% Art:		& {\FontTypewriter XXX}		& QU 	& {\FontTypewriter tt.mm.jjj hh:mm}	& BO 	& {\FontTypewriter hh:mm} 
% \\
% ZA			& {\FontTypewriter hh:mm}	& AO	& {\FontTypewriter hh:mm}			& EB 	& {\FontTypewriter hh:mm}
% \\\midrule
% Name: 		& {\FontTypewriter MUSTERMANN, Max}
% &
% Alter:		& {\FontTypewriter 32}
% &
% Geschlecht:	& {\FontTypewriter m}
% \\\addlinespace\midrule
% \multicolumn{6}{c}{\TabSub{Diagnosen}}
% \\
% \multicolumn{6}{p{.98\linewidth}}{{\FontTypewriter Diagnosen}}
% \\\midrule
% \multicolumn{6}{c}{\TabSub{Befunde}}
% \\
% Vitalwerte			& {\FontTypewriter 09:30} & RR - HF	& {\FontTypewriter 130/80}\MmHg* - {\FontTypewriter 100}\PerMin*	& Sp\ce{O2}	& {\FontTypewriter 99}\PC* 
% \\
% Vitalwerte			& {\FontTypewriter 09:50} & RR - HF	& {\FontTypewriter 130/80}\MmHg* - {\FontTypewriter 100}\PerMin*	& Sp\ce{O2}	& {\FontTypewriter 99}\PC* 
% \\
% BZ					& {\FontTypewriter xxx}\MgDl*& Temp.	& {\FontTypewriter xx}\,°C									& Schmerzen: 	&
% \\
% Weitere Befunde:	& \multicolumn{5}{p{.8\linewidth}}{{\FontTypewriter Weitere Befunde   }}
% \\
% NACA:				&		
% & 
% GCS:				&
% \\\addlinespace\midrule
% \multicolumn{6}{c}{\TabSub{Anamnese, Bemerkungen}}
% \\
% \multicolumn{6}{p{.98\linewidth}}{{\FontTypewriter m}}
% \\\midrule
% \multicolumn{6}{c}{\TabSub{Maßnahmen}}
% \\
% \multicolumn{6}{p{.98\linewidth}}{{\FontTypewriter m}}
% \\\midrule
% \multicolumn{6}{c}{\TabSub{Ziel}}
% \\
% Zielabt.:		& {\FontTypewriter m}
% &
% Trpt. m. Arzt: 	& {\FontTypewriter j/n}
% &
% Zustand:		& {\FontTypewriter k.A.}
% \\\addlinespace\bottomrule\addlinespace
% \end{tabularx}
% }
% {Ein Einsatzprotokoll}
% {}




\section{\HeadingKeyword{Unfallchirurgische} Diagnosen}

Im folgenden sind häufige (lateinische) Diagnosen bzw.
Abkürzungen als Referenz angeführt, welche man häufig auf
unfallchirurgischen Arztbriefen und Ambulanzprotokollen
findet.
Eine ausführliche Aufstellung findet sich in \cite{LatDiag4}.

{\small\Z{\TableBaseModification
\begin{longtable}{p{0.25\linewidth}p{0.2\linewidth}p{0.45\linewidth}}
\toprule\addlinespace
\TabHeading{Bezeichnung} & \TabHeading{Abkz.} & \TabHeading{Beschreibung}
\\\addlinespace\midrule\addlinespace\endhead
\multicolumn{3}{c}{{\TabSub{Wund- und Verletzungsarten}}}
\\\addlinespace\midrule\addlinespace
\IX{Abruptio} (ossea) 			& \I{Abrupt. oss.} 		&	Knöcherner Ausriß-Abriß
\\
\IX{Abscessus} 					& \I{Abscess.} 			&	Abszeß
\\
\IX{Amputatio traumatica} 		& \I{Amput. traumat.} 	&	Unfallbedingte Abtrennung
\\
\IX{Articulatio aperta} 		& \I{Artic. apert.} 		&	Gelenkseröffnung
\\\addlinespace
\IX{Bulla} haemorrhagica 		& \I[Bulla!haemorrh.]{Bulla haemorrh.}	&	Blutblase
\\
Bulla infecta 					& \I[Bulla!infect.]{Bulla infect.} 	&	infizierte Blase
\\
\IX{Bursa} aperta 				& \I{Bursa apert.} 		&	Eröffnung des Schleimbeutels
\\
\IX{Bursitis}  					& 					&	Schleimbeutelentzündung
\\
Bursitis purulenta  			& 					&	eitrige ~
\\\addlinespace
\IX{Cauterisatio}   			& \I{Cauteris.} 		&	Verätzung
\\
\IX{Combustio} (gradus \ldots)  & \I[Comb.]{Comb. (grad. \ldots)} &	Verbrennung {\ldots} Grades
\\
\IX{Congelatio}  				& \I{Congel.} 			&	Erfrierung
\\
\IX{Contusio}  					& \I{Cont.} 			&	Prellung, Quetschung
\\
\IX{Corpus alienum}  			& \I{Corp. alien.} 		&	Fremdkörper
\\
- ligneum, - ferreum, - vitreum 	&  					&	Holz-, Eisen-, Glas-Fremdkörper
\\
\IX{Corpus liberum} 			& 					&	Freier Gelenkskörper
\\\addlinespace
\IX{Decollement}  				& 					&	Hautabscherung
\\
\IX{Defectus cutis}  			& \I{Defect. cutis} 	&	Hautverlust
\\
\IX{Discissio} 					& \I{Disciss.} 			&	Durchtrennung
\\
\IX{Discissio nervi}  			& 					&	Burchtrennung des Nervs ( an den Fingern numeriert von 1 – 10 beginnend von der radialen Daumenseite )
\\
\IX{Distorsio}  					& \I{Dist.} 			&	Zerrung
\\\addlinespace
\IX{Empyema articulare}  		& 					&	Gelenkseiterung
\\
\IX{Enucleatio traumatica}  		& 					&	Auslösung im Gelenk
\\
\IX{Epidermiolysis}  			& \I{Epidermiol.} 		&	Ablederung der oberen Haut
\\
\IX{Epiphysiolysis}  			& \I{Epiphys.} 			&	Wachstumsfugenlösung
\\
\IX{Excoriatio}  				& \I{Excor.} 			&	Hautabschürfung
\\\addlinespace
\IX{Fissura} 					& \I{Fiss.} 			&	Sprung des Knochens
\\
\IX{Fistula} 					& 					&	Fistel
\\
\IX{Fractura}  					& \I{Fract.} 			&	Bruch des Knochens
\\
\IX[Fractura!aperta]{Fractura aperta}  			& \I[Fract.!apert.]{Fract. apert.} 	&	offener Bruch
\\
Fractura comminuta  		& 					&	Trümmerbruch
\\
Fractura duplex/triplex  	& 					&	zwei/dreifacher Bruch
\\
Fractura epiphysis  		& 					&	Bruch der Epiphyse
\\
Fractura subperiostalis  	& 					&	Grünholzbruch
\\\addlinespace
\IX{Hämarthros}  				& 					&	Bluterguß im Gelenk
\\
\IX{Hämatoma}  					& 					&	Bluterguß
\\
\IX{Hydrops}  					& 					&	Erguß im Gelenk
\\\addlinespace
Inflammatio  				& 					&	Hautrötung,Entzündung
\\\addlinespace
\IX{Laesio}  					& \I{Laes.} 		&	Verletzung, Läsion
\\\addlinespace
Laesio cartilaginis  		& \I[Laes.!cartilag.]{Laes. cartilag.} 	&	Verletzung des Knorpels
\\
Laesio nervi   				& \I[Laes.!nerv.]{Laes. nerv.} 		&	Verletzung des Nervs
\\
\IX{Luxatio}  					& \I{Lux.} 				&	Verrenkung
\\
Luxatio et fractura  		& \I{Luxfract.} 		&	Verrenkungsbruch
\\
Lyphangitis-adenitis  		& 					&	Lymphbahn-knotenentzündung
\\\addlinespace
\IX{Necrosis} cutis  			& 					&	abgestorbene Haut
\\\addlinespace
\IX{Panaritium} subcutaneum  	& 					&	Eiterung des Unterhautfettgewebes an Fingern oder Zehen
\\
Panaritium tendineum  		& 					&	eitrige Sehnenscheidenentz.
\\
\IX{Paralysis}   & 				&	Lähmung
\\
\IX{Paresis}   		& 					&	Teillähmung
\\
\IX{Paronychia}  				& 					&	Nagelbettentzündung
\\\addlinespace
\IX{Phlegmona}  					& 					&	Phlegmone
\\
\IX{Refractura}  				& \I{Refract.} 			&	Neuerlicher Bruch im Bereich eines alten Bruches
\\
\IX{Ruptura}  					& \I{Rupt.} 			&	Riß, Zerreißung
\\
\IX{Ruptura syndesmosis} 		& 					&	Syndesmosenzerreißung
\\\addlinespace
\IX{Seroma}  					& 					&	Erguß seröser Flüssigkeit
\\
\IX{Subluxatio}  				& \I{Sublux.} 			&	Teilverrenkung
\\
\IX{supraarticularis}  			& \I{supraartic.} 		&	über dem Gelenk
\\\addlinespace
\IX{Vulnus}		  				& \I{Vuln.} 			&	Wunde
\\\addlinespace
\IX[Vulnus!ictum]{Vulnus ictum}  				& \I[Vuln.!ict.]{Vuln. ict.} 			&	Stichwunde
\\
\II[Vulnus!lacero-contusum]{Vulnus lacero-contusum}  	& \I{Vlc.} 				&	Rißquetschwunde
\\
\IX[Vulnus!morsum!animalis]{Vulnus morsum animalis} 	& \I[Vuln.!mors. anim.]{Vuln. mors. anim.} 	&	Tierbiß
\\
\IX[Vulnus!morsum!hominis]{Vulnus morsum hominis}  		& \I[Vuln.!mors. hom.]{Vuln. mors. hom.}					&	Menschenbiß
\\
\IX[Vulnus!perforans]{Vulnus perforans}			& 					&	Durchdringende Stichwunde
\\
\IX[Vulnus!scissum]{Vulnus scissum} 				& \I[Vuln.!sciss.]{Vuln. sciss.} 		&	Schnittwunde
\\
\IX[Vulnus!sclopetarium]{Vulnus sclopetarium}			& \I[Vuln.!sclopet.]{Vuln. sclopet.}		&	Schußwunde
\\\addlinespace\midrule\addlinespace
\multicolumn{3}{c}{{\TabSub{Zusätze}}}
\\\addlinespace\midrule\addlinespace
\IX{recens}  		& \I{rec.} 		&	frisch (nur bei Unterscheidung gegen ältere Verletzungen anführen)
\\
\IX{non recens}  	& \I{non rec.} 		&	nicht frisch; bei bereits sichtbar infizierten oder über 8 Stunden alten Wunden oder offenen Brüchen, bei geschlossenen Brüchen nach 24 Stunden 
\\
\IX{infectus}, -a, -um	& 				&	infiziert
\\
\IX{inveteratus}  	& \I{invet.} 		&	veraltet
\\
\IX{apertus}  		& \I{apert.} 		&	offen
\\
\IX{ossea sanata}  	& \I{oss. san.} 	&	knöchern geheilt
\\
\IX{repositus}  		& \I{repos.} 		&	eingerichtet, reponiert
\\
\IX{utriusque}  		& \I{utr.} 			&	rechts und links
\\
\IX{bilateralis}  	& \I{bilat.} 		&	beiderseitig an einer Extremität (d.\,h. innen und aussen, radial und ulnar etc.)
\\
\IX{operatus}  		& \I{operat.} 		&	operiert
\\
\IX{suturatus}  		& \I{suturat.} 		&	genäht
\\
\IX{fixatus}  		& \I{fixat.} 		&	ruhiggestellt im Gips
\\
\IX{signum interrogationis} 	 & s.\,i. & fraglich, Zusatz bei unsicherer Diagnose
\\\addlinespace\midrule\addlinespace
\multicolumn{3}{c}{{\TabSub{Richtungs- und Seitenangaben}}}
\\\addlinespace\midrule\addlinespace
\IX{dexter}, -tra, -trum 		& \I{dext.} & rechts
\\
\IX{sinister}, -tra, -trum		& \I{sin.}	& links
\\
\IX{utriusque}					& \I{utr.} & Rechts und links
\\
\IX{bilateralis}, -e				& \I{bilat.} & beidseitig 
\\
\IX{medialis}, -e				& \I{med.} & innenseitig
\\
\IX{lateralis}, -e				& \I{lat.} & außenseitig
\\
\IX{dorsalis}, -e				& \I{dors.} & streckseitig
\\
\IX{palmaris}, -e				& \I{palm.} & beugeseitig an der Hand, hohlhandseitig
\\
\IX{plantaris}, -e				& \I{plant.} & beugeseitig am Fuß
\\
\IX{ulnaris}, -e					& \I{uln.} & innenseitig
\\
\IX{radialis}, -e				& \I{rad.} & außenseitig
\\
\IX{proximalis}, -e				& \I{prox.} & körpernahe
\\
\IX{distalis}, -e				& \I{dist.} & körperfern
\\
\IX{cranialis}, -e				& \I{cran.} & kopfwärts
\\
\IX{caudalis}, -e				& \I{caud.} & fußwärts
\\\addlinespace\bottomrule\addlinespace
\end{longtable}
}
}

\begin{Literary}
\fullcite{LatDiag4} \cite{LatDiag4}
\end{Literary}

\ChapterEnd